% res_refly_1m.tex -- DRAFT 23 Oct 2026 (thesis result R): every 0.55 m twin table re-flown on 1 m cables.
% Updated 23 Oct afternoon with batch 16 (the five 1 m worlds' tables) and reviewer 2's corrections.
% Sources: card docs/experiments/2026-10-23_refly_1m.md (bars frozen 23 Oct 02:15: each original card's own claim and
% thresholds); docs/analysis/refly_1m_table_2026-10-23.md (tools/refly_score.py, tools/thesis_runsets.py; batch 16 with
% the original cards' own scorers through a scratchpad wrapper); registry R2483-R2577 (1 m, lab PC, 4 slots; W4
% R2483-R2485, Tier A R2486-R2489/R2495-R2498/R2504-R2507, T15 R2490/R2491/R2499/R2500, W5 R2492/R2501, drop
% R2493/R2494/R2502/R2503, T7t R2508-R2537, appendix R2538-R2577); attribution re-flies R2580-R2583 on the same tree at
% 0.55 m; batch 16 R2604-R2628 (Tier B R2604-R2614, the push rows R2615-R2624, gain R2625-R2628; tree be34f95 + patch
% cf1ab9ad, whose only src/ change since R2483-R2577 is the rejoin's fixes, off in every R config).
% Reviewer [weak], 23 Oct 05:15, corrections applied here: C-L2 out of the table (both 1 m L-S runs R2564/R2577 miss S);
% C-D5's mapneg cells out on S and unresolved (footnote c); T7t orbit_v300 attributed by re-fly, not by argument; the
% read-back of R2485, R2521-R2536, R2542-R2550, R2564-R2577 lacks /tracker_2 (registry correction rows).
% Reviewer 2 [weak], 23 Oct 11:47 (Tier B, gain, pushes; every number reproduces to 0.01, all 25 runs valid, no bar
% re-set), corrections applied here: B4 a MARGINAL hold (R2606 0.19 deg under the line, R2612 2.50 under; n = 2, under the
% card's 0.5 deg tilt floor; the tier_b must-not-say stands: nothing says a gap up to 150 deg stays under the line with
% the planned pull); B5 NOT CONFIRMED, not "does not hold" (its falsifier, an abort or fault latch, did not fire); G165's
% drop (R2607/R2613) is an S MISS in both runs (the frozen S exempts only D1/D4/A2/A6; the ">=" reading came after the
% runs); C-D1 vertically and C-D4 EMPTY holds (D-Z0 without the integral already reads 0.06/0.06 cm).
% Tally (refly table doc, 11:47): 19 hold (W4, W5, T3, T8, T9, T15, drop, C-X1, C-X2, C-R1-R4, B1-B4, C-D1, C-D4), 5 do
% not (T7t new; C-L1, C-D2, C-D3, C-D5 as at 0.55 m), 2 unresolved (C-G1, C-G2), 1 not confirmed (B5); 0 void, 0
% skipped. S misses: mapneg R2542/R2549/R2543/R2550 (as at 0.55 m) and G165 R2607/R2613.
% tab:faults' "tilt over 30 deg" row: at 1 m it cites G165 (the integrator's call under Wesley's re-fly-at-1-m decision,
% SESSION_STATE 23 Oct 11:33; he can keep the 0.55 m history instead). tab:faults itself (Appendices.tex) is NOT edited
% here: it still cites the 0.55 m G150p drop.
% NOT REPORTED (Wesley 23 Oct, decisions.md DECIDED 2026-10-23 (9)): the straight-descent landing L-S (C-L2) at 1 m. Its
% two runs stalled 0.14-0.17 m up on the rigid 1 m rods and the stall exit fired; at 0.55 m on this tree it passes
% (R2582/R2583), so the stall is the twin's rigid 1 m rod. (The decision calls it "the 2-quad landing"; the config,
% app_land_straight, has three carriers.)
% The tree changed since the 0.55 m runs (card section 3): land_rearm_touchdown on in the twin (17 Oct); cable_accel_cap
% 12 -> 14.5 m/s^2 (6.6 -> 8 N, 21 Oct); the twin's mocap velocity filter none -> ema (22 Oct); the 1 m defaults (the
% freed UAV's path-clear distance 1.225 m); the 5 s hover hold before every path (T7t's 15 cells, X0-O0/O1). T7t's
% figure-eight is size 0.5 (as flown at 0.55 m; one variable), the rig's is 0.75.
% Must not say (card section 10): no 1 m twin number against a 0.55 m rig flight; "the cable changed X" only with an
% attribution re-fly; W4's 0.55 m column is n = 1; the integral gains are the 0.55 m ones, not re-tuned; the rod error is
% +5 cm absolute (5 % of the rod at 1 m, 9 % at 0.55 m). D-X0's larger steady offset at 1 m (6.6 against 4.8 cm) is not
% put on the cable (no attribution re-fly).

\subsection{The claims on \texorpdfstring{1\,m}{1 m} cables}
\label{sec:res_refly}

This section asks whether the claims that the twin tables of Sections~\ref{sec:res_carry} and~\ref{sec:res_leave}
make on 0.55\,m rods still hold on 1\,m cables, the cables the rig flies from 28 October. Every table was re-flown at
1\,m with its own arms, layouts and repeats, and each claim was scored against its own card's bars, unchanged and
frozen before the re-fly. The baseline is the set of 0.55\,m runs that each table cites. Of the 27 claims scored, 19
hold, 5 do not, 2 stay unresolved as they were at 0.55\,m, and 1 is not confirmed (Tables~\ref{tab:refly_claims},
\ref{tab:refly_steady}, \ref{tab:refly_push} and~\ref{tab:refly_margin}). One of the holds is marginal, and the vertical part of two is
empty: the vertical push left nothing for the remedy to remove.

The comparison is not one variable. Between the 0.55\,m runs and the re-fly the code tree moved: the touchdown test
is re-armed at ring touchdown, the cap on the rod pull that the trackers feed forward was raised, the twin's motion
capture filters velocities as the rig's does, and every path now starts after a 5\,s hover hold. A difference between
the two columns is therefore put on the cable only where the 0.55\,m configuration, re-flown on the new tree,
separates the two (Table~\ref{tab:refly_attrib}). The integral's gains are the ones identified on 0.55\,m rods; they
were not re-tuned for 1\,m. The tables of the release margin, the integral's gain and the constant pushes needed 1\,m
worlds of their own; these were built from the rig's geometry with the cable the only change and flown after the rest,
on the same code but for the rejoin's fixes (Section~\ref{sec:res_rejoin}), which these runs leave off.

\begin{table}[htbp]
  \centering
  \footnotesize
  \setlength{\tabcolsep}{3pt}
  \caption[The twin's claims on 1\,m cables: leaving and carrying]{The release and carry claims of the twin's
  0.55\,m tables, re-flown on 1\,m cables and scored against each card's own bars, frozen before the runs: four UAVs in
  the release rows and three in the carry row, the 0.86\,kg ring at 1.0\,m, floor start, the default controller; runs
  per cell in brackets. The 1\,m runs were flown on one code tree, newer than that of the 0.55\,m runs (see text).
  Release rows: the tilt rise $\Delta\theta$ and the resolution rule of Section~\ref{sec:metrics} (the gap between the
  two arms' means, and the gap the rule needs). $^a$\,With detection off the 1\,m ring peaked lower than the 0.55\,m
  ring, and no 1\,m run ended in an abort in its LAND (0.55\,m: three of three).}
  \label{tab:refly_claims}
  \begin{tabularx}{\textwidth}{>{\raggedright\arraybackslash}p{4.9cm}>{\raggedright\arraybackslash}X>{\raggedright\arraybackslash}X>{\raggedright\arraybackslash}p{1.8cm}}
    \toprule
    Claim (table) & 0.55\,m & 1\,m & At 1\,m \\
    \midrule
    \multicolumn{4}{l}{\emph{Leaving the team}} \\
    The freed UAV is disarmed alone, hover (Table~\ref{tab:T13}) & rest to disarm 0.28\,s; $\Delta\theta$ 6.1$^\circ$ (1)
      & 0.35--0.36\,s; 7.0--7.2$^\circ$ (3) & holds \\ % R1887 | R2483-R2485 (P3 <= 0.003 m, P4 0)
    The freed UAV is disarmed alone, circle (Table~\ref{tab:T13}) & 0.27--0.28\,s; 6.7--7.1$^\circ$ (2)
      & 0.32--0.34\,s; 7.9--8.0$^\circ$ (2) & holds \\ % R1913/R1914 | R2492/R2501
    Pull source, 120$^\circ$ gap: measured $<$ planned (Table~\ref{tab:T9}) & 6.4 against 16.8$^\circ$; gap 10.4, needed 4.4 (3)
      & 7.3 against 17.1$^\circ$; gap 9.8, needed 1.2 (3) & holds \\ % T3: A1 vs A2
    Detection, planner's pull: on $<$ off (Table~\ref{tab:T9}) & 16.8 against 28.2$^\circ$; gap 11.4, needed 1.0 (3)
      & 17.1 against 21.7$^\circ$$^a$; gap 4.6, needed 1.0 (3) & holds \\ % T8: A2 vs A6
    Detection, measured pull: no sign (Table~\ref{tab:T9}) & gap $-$0.8, needed 4.4: unresolved (3)
      & gap $+$0.3, needed 1.4: unresolved (3) & holds \\ % T9: A1 vs A5
    Pull source, 60$^\circ$ circle, no event: planned $<$ measured (Table~\ref{tab:T9}) & 0.53 against 0.93\,cm; gap 0.39, needed 0.30 (2)
      & 0.61 against 0.96\,cm; gap 0.35, needed 0.30 (2) & holds \\ % T15: A8 vs A7
    A second loss, or a 240$^\circ$ gap, drops the ring and lands every UAV (Table~\ref{tab:faults})
      & 4 of 4: dropped 0.22--0.23\,s after the loss; UAV tilt at most 13.6$^\circ$; all down in 8.5--8.9\,s
      & 4 of 4: 0.20--0.22\,s; at most 13.7$^\circ$; 8.1--8.5\,s & holds \\ % D1 R1915/R1916, D4 R1917/R1918 | R2493/R2502, R2494/R2503
    \midrule
    \multicolumn{4}{l}{\emph{Carrying}} \\
    Every cell inside the envelope at 0.55\,m stays inside (Table~\ref{tab:refly_T7t})
      & 14 of 15 cells inside; the spin at 0.3\,m/s outside & 14 of 15 inside; the orbit at 0.3\,m/s outside, the spin
      at 0.3\,m/s inside & does not hold \\
    \bottomrule
  \end{tabularx}
\end{table}

\begin{table}[htbp]
  \centering
  \footnotesize
  \setlength{\tabcolsep}{3pt}
  \caption[The twin's claims on 1\,m cables: steady errors]{The steady-error claims of the twin's 0.55\,m tables,
  re-flown on 1\,m cables as Table~\ref{tab:refly_claims}: three carriers, hover unless stated. $^b$\,A UAV more than
  10\,cm above its rest when the fleet was declared landed; in every case, at both lengths, one UAV was stronger than its
  thrust map; the 1\,m count includes the landings of Table~\ref{tab:refly_push}, all clean. $^c$\,These runs also
  disarm a UAV in the air (footnote~$b$).}
  \label{tab:refly_steady}
  \begin{tabularx}{\textwidth}{>{\raggedright\arraybackslash}p{4.9cm}>{\raggedright\arraybackslash}X>{\raggedright\arraybackslash}X>{\raggedright\arraybackslash}p{1.8cm}}
    \toprule
    Claim (table) & 0.55\,m & 1\,m & At 1\,m \\
    \midrule
    The tracker's relaxed initial condition ($\epsilon = 0.025$) raises the hover by at least 2\,cm (Table~\ref{tab:x0}) & $+$13.1\,cm (2) & $+$18.7\,cm (2) & holds \\ % C-X1
    Pinning the tracker's initial condition ($\epsilon = 0$) at least halves the orbit's error (Table~\ref{tab:x0}) & 0.74 against 3.87\,cm, ratio 0.19 (2)
      & 0.63 against 3.98\,cm, ratio 0.16 (2) & holds \\ % C-X2
    Inside the integral's gate every model error is removed to within 1.0\,cm (Table~\ref{tab:robust})
      & worst 0.87\,cm & worst 0.41\,cm & holds \\ % C-R1
    With the mass typed 20\,\% high, or the rods typed 5\,cm long and not measured, the integral changes the error by
      at most 1.0\,cm (Table~\ref{tab:robust}) & 0.01 and 0.00\,cm; the unmeasured rods leave $+$17.1\,cm
      & 0.01 and 0.00\,cm; $+$17.8\,cm & holds \\ % C-R2: rod+5nomeas +17.09 -> +17.75 both columns
    Measuring the rods absorbs a 5\,cm length error (Table~\ref{tab:robust}) & within 0.00\,cm of the error-free hover
      & within 0.00\,cm & holds \\ % C-R3
    Every robustness flight lands, no UAV past 45$^\circ$ (Table~\ref{tab:robust}) & 28 of 28 & 28 of 28 & holds \\ % C-R4
    The default landing tips no UAV and disarms none in the air$^b$ (Table~\ref{tab:land})
      & no tip; 4 of 82 disarmed in the air, 0.26--0.53\,m up & no tip; 4 of 82 disarmed in the air, 0.17\,m up;
      outward at most 0.8\,cm, tilt at most 5.3$^\circ$ & does not hold, as at 0.55\,m \\ % C-L1: R2542/R2549/R2543/R2550; 68 + the 14 push-arm landings R2615-R2628, all clean
    The model integral removes a thrust-map error's horizontal error better than the target shift
      (Table~\ref{tab:disturbance}) & 0.01 against 0.01\,cm: unresolved (2) & 0.00 against 0.01\,cm: unresolved$^c$ (2)
      & does not hold, as at 0.55\,m \\ % C-D5: D-MM vs D-MR
    \bottomrule
  \end{tabularx}
\end{table}

\begin{table}[htbp]
  \centering
  \footnotesize
  \setlength{\tabcolsep}{3pt}
  \caption[The twin's claims on 1\,m cables: pushes and the integral's gain]{The claims of the twin's 0.55\,m tables on
  a constant push and on the integral's gain, re-flown on 1\,m cables as Table~\ref{tab:refly_claims}, on the 1\,m world
  built for them: three carriers in hover, a 0.6\,N vertical or 0.4\,N horizontal push on the ring for 15\,s. Push
  rows: the steady offset (vertical; horizontal) and the force the integral learned. Gain rows: the peak or the
  overshoot of the vertical push against the default's. $^a$\,Without any integral the vertical push already leaves
  almost no offset (0.06\,cm at 1\,m, at most 0.03\,cm at 0.55\,m): the measured pull takes it up
  (Section~\ref{sec:res_force}), so these vertical holds show nothing.}
  \label{tab:refly_push}
  \begin{tabularx}{\textwidth}{>{\raggedright\arraybackslash}p{4.9cm}>{\raggedright\arraybackslash}X>{\raggedright\arraybackslash}X>{\raggedright\arraybackslash}p{1.8cm}}
    \toprule
    Claim (table) & 0.55\,m & 1\,m & At 1\,m \\
    \midrule
    The model integral removes a constant push to within 1.0\,cm (Table~\ref{tab:disturbance})
      & 0.01--0.02; $-$0.01--0.00\,cm (2) & 0.03--0.06; 0.05\,cm (2) & holds; vertically empty$^a$ \\ % C-D1: D-ZM R2617/R2622, D-XM R2619/R2624
    Without it the push leaves at least 2\,cm (Table~\ref{tab:disturbance})
      & 0.01--0.03; 4.81--4.82\,cm (2) & 0.06; 6.56--6.58\,cm (2) & does not hold vertically, as at 0.55\,m \\ % C-D2: D-Z0 R2615/R2620, D-X0 R2618/R2623
    The integral learns the applied force: 0.45--0.75\,N vertically, 0.30--0.50\,N horizontally (Table~\ref{tab:disturbance})
      & 0.099--0.102; 0.325\,N (2) & 0.112--0.120; 0.325--0.327\,N (2) & does not hold vertically, as at 0.55\,m \\ % C-D3
    The target shift removes the vertical push to within 1.0\,cm (Table~\ref{tab:disturbance})
      & 0.00--0.02\,cm (2) & $-$0.07 to $-$0.06\,cm (2) & holds; empty$^a$ \\ % C-D4: D-ZR R2616/R2621
    Doubling the integral's gain lowers the vertical push's peak (Table~\ref{tab:gain})
      & 1.26--1.44 against 1.25--1.41\,cm: unresolved (2) & 1.42--1.44 against 1.42--1.47\,cm: unresolved (2)
      & unresolved, as at 0.55\,m \\ % C-G1: G-K8 R2625/R2627 vs D-ZM R2617/R2622
    Keeping the speed gate only for 5\,s after the lift lowers the overshoot (Table~\ref{tab:gain})
      & 1.31--1.33 against 1.39\,cm: under the floor (2) & 1.39--1.40 against 1.41--1.44\,cm: under the floor (2)
      & unresolved, as at 0.55\,m \\ % C-G2: G-S5 R2626/R2628 vs D-ZM
    \bottomrule
  \end{tabularx}
\end{table}

The release claims keep their signs at 1\,m. The freed UAV is still disarmed alone, the measured pull still halves the
tilt rise or better at a 120$^\circ$ gap, and a second loss still drops the ring and lands every UAV. With detection
off and the planner's pull, the 1\,m ring rose less than the 0.55\,m ring did and no run ended in an abort, so the
claim holds on a smaller gap. The steady-error claims hold with the same signs: the tracker's relaxed initial condition raises
the hover further on the longer cables, and the unmeasured 5\,cm rod error left the same height error, although the same 5\,cm
is a smaller part of a 1\,m rod. The tables reported without a bar (the hover and paths of Table~\ref{tab:traj}, the
flight of Figure~\ref{fig:flight}) were flown at 1\,m as well, and the landing-lean rows of Table~\ref{tab:T9} were not
re-scored. [Table~\ref{tab:traj} and Figures~\ref{fig:flight}, \ref{fig:traj} and~\ref{fig:detach}: to regenerate from
the 1\,m runs.]
% traj: R2544/R2551 hover 0.03/0.04 cm mean, 0.06 deg (0.55 m 0.04/0.04, 0.06); flight: R2544 (S PASS); K1 not scored.

Four of the five claims that do not hold at 1\,m failed in the same way at 0.55\,m. The landing still disarms the fleet
with a UAV in the air when that UAV is stronger than its thrust map (Section~\ref{sec:sim_landing}); at 1\,m that UAV
was lower when it happened, and the landings of the push rows were all clean. The model integral and the target shift
still leave a thrust-map error's horizontal error equally small, so the integral's advantage there remains unshown.
The vertical push still leaves no offset to remove, so the claims that the push shows without the integral and that
the integral learns its force fail vertically, as they did, and hold horizontally, where the integral removes the
offset and learns the force it learned at 0.55\,m. The integral's gain and its speed gate still make no difference
the resolution rule can resolve. The fifth, the carry envelope, is new at 1\,m.

\subsubsection{The carry envelope at \texorpdfstring{1\,m}{1 m}}

Table~\ref{tab:refly_T7t} gives the carry envelope of Table~\ref{tab:T7t} at both lengths, scored on the clock-free
path error: the rms, over the path, of the ring's distance across the path it was sent and of its height error, with
no schedule in it. The bars are the carry card's: a path error of at most 25\,mm, a peak ring tilt over the path of at
most 6$^\circ$ and a radius ratio of at least 0.95. At 1\,m the orbit at 0.3\,m/s leaves the envelope on two bars at
once, by a small margin on each, and the spin at 0.3\,m/s, outside at 0.55\,m on one run's tilt, comes inside.
Across the envelope the path error rose in every cell but the line's, the radius ratio fell in every cell, and the
path tilt fell or stayed.

\begin{table}[htbp]
\centering
\footnotesize
\setlength{\tabcolsep}{2.6pt}
\caption[The carry envelope on 1\,m cables]{The three-carrier carry envelope in the rig twin on 0.55\,m rods and on
1\,m cables: the 0.86\,kg ring at 1.0\,m, floor start, the default controller, two runs per cell on one code tree per
length (the 1\,m tree is the newer); each cell gives the range over its two runs. Path error: the clock-free path
error (bar 25\,mm); path tilt: the peak ring tilt over the path (bar 6$^\circ$); ratio: the ring's path radius over the
reference's (for the line, the ratio of amplitudes about its centre; bar 0.95). Speed: the peak path speed. Beside
them, not scored: the frozen ring RMSE of Table~\ref{tab:T7t} at 1\,m (its frozen peak tilt was at most 3.0$^\circ$ in
every 1\,m run). $^a$\,Over a bar: the cell is
outside the envelope. $^b$\,Four decimals near the bar.}
\label{tab:refly_T7t}
\begin{tabular}{llcccccc}
\toprule
 & & \multicolumn{2}{c}{Path error (mm)} & \multicolumn{2}{c}{Path tilt ($^\circ$)} & Ratio & Frozen RMSE \\
\cmidrule(lr){3-4}\cmidrule(lr){5-6}
Path & Speed (m/s) & 0.55\,m & 1\,m & 0.55\,m & 1\,m & 1\,m & 1\,m (mm) \\
\midrule
Circle & 0.125 & 2.5 & 3.4--3.5 & 0.86--1.11 & 0.69--0.74 & 0.995 & 9.2--9.8 \\
 & 0.2 & 3.7--6.3 & 7.9--8.1 & 1.45--1.58 & 1.01--1.13 & 0.987--0.988 & 11.7--11.9 \\
 & 0.3 & 11.0--11.7 & 14.7--15.1 & 2.04--2.79 & 1.55--1.60 & 0.975 & 14.9--15.1 \\
\addlinespace
Figure-eight & 0.125 & 3.8--6.8 & 8.1--8.2 & 0.97--1.06 & 0.87--0.91 & 0.987 & 12.0 \\
 & 0.2 & 9.3--10.0 & 14.3--15.0 & 1.48--1.73 & 1.40--1.47 & 0.972--0.974 & 15.5--16.3 \\
 & 0.3 & 15.4--17.6 & 20.5--20.8 & 2.31--2.32 & 2.26--2.33 & 0.955--0.957 & 19.5--19.8 \\
\addlinespace
Line & 0.125 & 1.2--1.4 & 1.1--1.2 & 0.91--0.99 & 0.77--0.79 & 0.990 & 7.8 \\
 & 0.2 & 1.7 & 1.5 & 1.54--1.88 & 1.03--1.18 & 0.973--0.974 & 11.3 \\
 & 0.3 & 2.4--2.6 & 1.5--1.6 & 2.42--2.96 & 2.14--2.18 & 0.951 & 18.0--18.1 \\
\addlinespace
Orbit & 0.125 & 4.1--5.2 & 5.2--5.5 & 0.98--1.07 & 0.75--0.79 & 0.990 & 8.5--8.7 \\
 & 0.2 & 9.0--9.9 & 13.4--13.7 & 1.67--1.88 & 1.09--1.13 & 0.974 & 14.0--14.3 \\
 & 0.3 & 13.2--15.2 & 25.7--25.9$^a$ & 2.49--2.85 & 1.76 & 0.9495--0.9499$^{a,b}$ & 23.8--24.0 \\
\addlinespace
Spin & 0.125 & 2.5--3.5 & 4.5 & 1.05--1.09 & 0.71--0.77 & 0.993--0.994 & 9.5--9.9 \\
 & 0.2 & 3.1--6.0 & 8.2--8.3 & 1.94--2.20 & 1.02--1.15 & 0.987--0.988 & 11.6--12.3 \\
 & 0.3 & 9.0--12.1 & 15.3--15.7 & 3.76--6.31$^a$ & 1.81--1.83 & 0.976 & 15.0--15.5 \\
\bottomrule
\end{tabular}
\end{table}
% 0.55 m: R1847-R1876 (freeze tree, flight code 4b3560e), re-scored on the 22 Oct ruler by tools/refly_score.py (each
% reproduces the carry card's verdicts). 1 m: R2508-R2537 (re-fly tree). R2511 (orbit_v125) is a runner stall 2.4 s
% after LANDED, the flight complete: scored valid (Wesley 23 Oct (4)). Five cells rose more than 1.5x (circle 0.2,
% figure-eight 0.125 and 0.2, orbit 0.3, spin 0.2): informational (card section 5, A), not attributed except orbit 0.3.

The attribution re-fly puts most of the orbit's rise on the code tree, not on the cable
(Table~\ref{tab:refly_attrib}): the same 0.55\,m configuration flown on the new tree already came within a few
millimetres of the bar, and the 1\,m cable added about 2.6\,mm and took the cell over it. Which of the tree's changes
raised the error is not separated. Figure~\ref{fig:refly_orbit} shows where along the orbit the ring runs inside its
path in the three conditions.

\begin{table}[htbp]
\centering
\footnotesize
\caption[Attribution of the orbit's rise]{Attribution of the orbit at 0.3\,m/s, the one cell that leaves the
envelope at 1\,m: the same configuration on 0.55\,m rods flown on the code tree of the 0.55\,m table and on the
re-fly's tree, and the 1\,m cell on the re-fly's tree; two runs each. Bars: path error at most 25\,mm, ratio at least
0.95.}
\label{tab:refly_attrib}
\begin{tabular}{lccc}
\toprule
Cables, code tree & Path error (mm) & Ratio & Path tilt ($^\circ$) \\
\midrule
0.55\,m, the 0.55\,m table's tree & 13.2--15.2 & 0.975--0.977 & 2.49--2.85 \\
0.55\,m, the re-fly's tree & 22.7--23.7 & 0.954--0.956 & 1.75--1.77 \\
1\,m, the re-fly's tree & 25.7--25.9 & 0.9495--0.9499 & 1.76 \\
\bottomrule
\end{tabular}
\end{table}
% R1871/R1872 | R2580/R2581 | R2521/R2536.

\pendingfig{fig:refly_orbit}{The orbit at 0.3\,m/s, the cell that leaves the carry envelope on 1\,m cables: plan view
of the ring's centre over one lap against the reference circle (dotted), on 0.55\,m rods with the 0.55\,m table's code
tree (grey), on 0.55\,m rods with the re-fly's tree (light gold) and on 1\,m cables with the re-fly's tree (gold);
(b) the signed distance across the path along the lap, positive inside, with the tightest quarter shaded.}{the
orbit runs at 0.3\,m/s of the three conditions of Table~\ref{tab:refly_attrib}}
% Data: R1871 (0.55 m, old tree), R2580 (0.55 m, re-fly tree), R2521 (1 m); run.csv payload_x/y and payload_ref_x/y,
% the path from metrics.path_through / project_on_path as tools/n2_bands.path_metrics_clock_free. Tool: a new function
% in tools/thesis_sim_figures.py through the thesis-figure skill.

\subsubsection{The release margin at \texorpdfstring{1\,m}{1 m}}

Table~\ref{tab:refly_margin} gives the margin curve of Table~\ref{tab:T9} on 1\,m cables: a hover release on the
rig's plates that leaves the survivors a gap of 120 to 180$^\circ$, and with it a smaller margin of the ring's centre
of mass inside them. The curve keeps its order: the smaller the margin, the larger the tilt rise, and the measured pull
still rises less than the planner's at 150$^\circ$. Its shape is cleaner at 1\,m, with every step between neighbouring
gaps now resolved with its sign.

Two things changed at the drop line. With the planner's pull at 150$^\circ$ the 1\,m ring no longer reached the
30$^\circ$ line, which it did in both runs at 0.55\,m; one run came within a fraction of a degree of it, so the claim
holds only by its letter, and nothing here says that a gap of up to 150$^\circ$ stays under the line on the planner's
pull. At 165$^\circ$ with the measured pull, under the line at 0.55\,m, the 1\,m ring crossed it in both runs: the
trackers' drop rule dropped the ring and every UAV landed. That is the exit of Table~\ref{tab:faults} working as
built, and at 1\,m these are the runs that stand for its tilt row. The re-fly's safety bar did not allow for a drop
there, so both runs miss it; and in both the freed UAV, already at rest, was disarmed with the fleet when the drop's
landing was declared, not alone. Nothing aborted and no fault fired, which is what that claim's falsifier asks, so it
is not confirmed rather than failed.

\begin{table}[htbp]
  \centering
  \footnotesize
  \setlength{\tabcolsep}{3pt}
  \caption[The release margin on 1\,m cables]{The release-margin claims of the twin's 0.55\,m tables, re-flown on
  1\,m cables and scored as Table~\ref{tab:refly_claims}: four UAVs on the rig's plates, an unannounced hover release
  that leaves the survivors a gap of 120 to 180$^\circ$, the measured pull unless stated; two runs per gap, three at
  120$^\circ$, one for the default rules at 180$^\circ$. Values: the tilt rise $\Delta\theta$ ($^\circ$), the means;
  $\geq$, cut at the drop or the latch. $^a$\,One run came 0.19$^\circ$ under the line, less than the re-fly's
  0.5$^\circ$ floor for a tilt, and its repeat 2.69$^\circ$ under it. $^b$\,Every UAV was disarmed at rest and nothing aborted;
  at 165$^\circ$ the freed UAV was disarmed with the fleet when the drop's landing was declared. $^c$\,A drop the
  re-fly's safety bar did not allow for: both runs miss it.}
  \label{tab:refly_margin}
  \begin{tabularx}{\textwidth}{>{\raggedright\arraybackslash}p{4.9cm}>{\raggedright\arraybackslash}X>{\raggedright\arraybackslash}X>{\raggedright\arraybackslash}p{1.8cm}}
    \toprule
    Claim (table) & 0.55\,m & 1\,m & At 1\,m \\
    \midrule
    A smaller margin gives a larger rise: gaps of 120, 150 and 180$^\circ$, both steps resolved (Table~\ref{tab:T9})
      & 6.4, 16.8 and $\geq$61.2; gaps 10.4 and 44.4, needed 4.4 and 2.0 & 7.3, 17.9 and $\geq$62.2; gaps 10.6 and 44.3,
      needed 1.2 and 1.2 & holds \\ % B1: G120 (A1) / G150 R2605/R2611 / G180 R2608/R2614
    The curve's shape: no step between neighbouring gaps resolved the other way (Table~\ref{tab:T9})
      & 120 to 135$^\circ$ unresolved (gap 2.0, needed 4.4); the rest resolved & every step resolved with its sign
      (120 to 135$^\circ$: gap 3.4, needed 1.2) & holds \\ % B2: G135 R2604/R2610 10.72, G165 R2607/R2613 >= 30.16
    Pull source, 150$^\circ$ gap: measured $<$ planned (Table~\ref{tab:T9}) & 16.8 against $\geq$29.5; gap 12.7, needed 4.3
      & 17.9 against 28.4; gap 10.5, needed 5.0 & holds \\ % B3: G150 vs G150p R2606/R2612
    Every absolute peak under the 30$^\circ$ drop line, gaps 120--150$^\circ$, either pull (Table~\ref{tab:T9})
      & the planner's pull at 150$^\circ$ reached the line in 2 of 2 & largest 29.8$^\circ$, the planner's pull at
      150$^\circ$; no drop$^a$ & holds, marginal \\ % B4: R2606 29.81, R2612 27.31
    Gaps 120--165$^\circ$, the planner's pull at 150$^\circ$ included: no abort, every UAV lands, the freed UAV is
      disarmed alone & 11 of 11; the freed UAV disarmed 0.19--0.43\,s after it came to rest
      & 9 of 11, 0.33--0.37\,s; at 165$^\circ$ disarmed with the fleet$^b$ & not confirmed \\ % B5: G165 0.04 s after LANDED
    \midrule
    \multicolumn{4}{l}{\emph{Reported, no bar (Table~\ref{tab:faults})}} \\
    Ring tilt over 30$^\circ$: the ring dropped, every UAV landed & the planner's pull at 150$^\circ$: 2 of 2
      & the measured pull at 165$^\circ$: 2 of 2, all down in 8.2\,s$^c$; the planner's pull at 150$^\circ$: no drop
      & -- \\ % G165 R2607/R2613 (trackers' 30 deg request +0.50/+0.51 s); 0.55 m G165 landed, abs peaks 28.10/27.84
    A 180$^\circ$ gap with the drop rule off: the fleet disarmed in the air at the 60$^\circ$ latch
      & 1.45--1.62\,m up (2) & 1.84--2.00\,m up (2) & -- \\ % G180 R2608/R2614 (fault at +0.64/+0.65 s)
    The same gap on the default rules: dropped by the gap rule at detection & $\geq$24.1 (1) & $\geq$23.2 (1)
      & -- \\ % G180d R2609 (+0.30 s; 0.55 m +0.32 s)
    \bottomrule
  \end{tabularx}
\end{table}
% Per run (refly table doc, "Batch 16"): G120 R2486/R2495/R2504 6.90/7.45/7.50; G135 10.42/11.02; G150 18.17/17.58;
% G150p 29.62/27.10 (abs 29.81/27.31; back in band 10.11/8.12 s); G165 >=30.49/>=29.84 (abs 30.30/29.57); G180
% >=62.46/>=61.91; G180d >=23.15. A agrees on every arm.

\paragraph{Answer.} On 1\,m cables the twin's claims hold as they did on 0.55\,m rods, with three changes of verdict:
the orbit at 0.3\,m/s now sits just outside the carry envelope, mostly from the code tree and a little from the cable;
the planner's pull at a 150$^\circ$ gap now stays under the drop line, but only just; and at a 165$^\circ$ gap, where
the 1\,m ring now crosses the line and is dropped, the freed UAV is disarmed with the fleet rather than alone. The
claims that failed at 0.55\,m fail the same way at 1\,m, and the gain's two stay unresolved. The 1\,m tables
therefore replace the 0.55\,m twin tables, with this section as the comparison; at 1\,m the drop at 165$^\circ$ stands
for the tilt row of Table~\ref{tab:faults}. The twin's rigid 1\,m rods remain its largest difference from the rig's
slack cables.
